%% TRANSCRIBED (PDF 37-48)

\apphead{Equipartition and \textit{L}-functions}%
  {EQUIPARTITION AND \textit{L}-FUNCTIONS}

\subsechead{A.1}{Equipartition}

Let $X$ be a compact topological space and $C(X)$ the Banach space of continuous, complex-valued functions on $X$, with its usual norm $\|f\| = \sup_{x \in X} |f(x)|$. For each $x \in X$ let $\delta_x$ be the Dirac measure associated to $x$; if $f \in C(X)$, we have $\delta_x(f) = f(x)$. Let $(x_n)_{n \geq 1}$ be a sequence of points of $X$. For $n \geq 1$, let $\mu_n = (\delta_{x_1} + \cdots + \delta_{x_n})/n$. Let $\mu$ be a Radon measure on $X$ (i.e. a continuous linear form on $C(X)$, cf. Bourbaki, Int., chap. III, \S 1). The sequence $(x_n)$ is said to be $\mu$-\emph{equidistributed}, or $\mu$-\emph{uniformly distributed}, if $\mu_n \longrightarrow \mu$ weakly as $n \to \infty$, i.e. if $\mu_n(f) \longrightarrow \mu(f)$ as $n \to \infty$ for any $f \in C(X)$. Note that this implies that $\mu$ is positive and of total mass $1$. Note also that $\mu_n(f) \longrightarrow \mu(f)$ means that
\[
\mu(f) = \lim_{n \to \infty} \frac{1}{n} \sum_{i=1}^{n} f(x_i).
\]

\begin{lemma}[1]\label{lem:I.A.1.1}
Let $(\phi_a)_a$ be a family of continuous functions on $X$ with the property that their linear combinations are dense in $C(X)$. Suppose that, for all $a$, the sequence $(\mu_n(\phi_a))_{n \geq 1}$ has a limit. Then the sequence $(x_n)$ is equidistributed with respect to some measure $\mu$; it is the unique measure such that $\mu(\phi_a) = \lim_{n \to \infty} \mu_n(\phi_a)$ for all $a$.

If $f \in C(X)$, an argument using equicontinuity shows that the sequence $(\mu_n(f))$ has a limit $\mu(f)$, which is continuous and linear in $f$; hence the lemma.
\end{lemma}

\begin{proposition}[1]\label{prop:I.A.1.1}
Suppose that $(x_n)$ is $\mu$-equidistributed. Let $U$ be a subset of $X$ whose boundary has $\mu$-measure zero, and, for all $n$, let $n_U$ be the number of $m \leq n$ such that $x_m \in U$. Then $\lim_{n \to \infty} (n_U/n) = \mu(U)$.
\end{proposition}

\begin{remarks}
\begin{nlist}
\item Let $U^\circ$ be the interior of $U$. We have $\mu(U^\circ) = \mu(U)$. Let $\varepsilon > 0$. By the definition of $\mu(U^\circ)$ there is a continuous function $\phi \in C(X)$, $0 \leq \phi \leq 1$, with $\phi = 0$ on $X - U^\circ$ and $\mu(\phi) \geq \mu(U) - \varepsilon$. Since $\mu_n(\phi) \leq n_U/n$ we have
\[
\liminf_{n \to \infty} n_U/n \geq \lim_{n \to \infty} \mu_n(\phi) = \mu(\phi) \geq \mu(U) - \varepsilon,
\]
from which we obtain $\liminf n_U/n \geq \mu(U)$. The same argument applied to $X - U$ shows that $\liminf_{n} (n - n_U)/n \geq \mu(X - U)$. Hence $\limsup n_U/n \leq \mu(U) \leq \liminf n_U/n$, which implies the proposition.

\item Let $X = [0,1]$, and let $\mu$ be the Lebesgue measure. A sequence $(x_n)$ of points of $X$ is $\mu$-equidistributed if and only if for each interval $[a,b]$, of length $d > 0$ in $[0,1]$, the number of $m \leq n$ such that $x_m \in [a,b]$ is equivalent to $dn$ as $n \to \infty$.

\item Let $G$ be a compact group and let $X$ be the space of conjugacy classes of $G$ (i.e. the quotient space of $G$ by the equivalence relation induced by inner automorphisms of $G$). Let $\mu$ be a measure on $G$; its image of $G \to X$ is a measure on $X$, which we also denote by $\mu$. We then have the following.
\end{nlist}
\end{remarks}

\begin{proposition}[2]\label{prop:I.A.1.2}
The sequence $(x_n)$ of elements of $X$ is $\mu$-equidistributed if and only if for any irreducible character $\chi$ of $G$ we have
\[
\lim_{n \to \infty} \frac{1}{n} \sum_{i=1}^{n} \chi(x_i) = \mu(\chi).
\]
\end{proposition}

The map $C(X) \to C(G)$ is an isomorphism of $C(X)$ onto the space of central functions on $G$; by the Peter-Weyl theorem, the irreducible characters $\chi$ of $G$ generate a dense subspace of $C(X)$. Hence the proposition follows from Lemma~\ref{lem:I.A.1.1}.

\begin{corollary}[1]\label{cor:I.A.1.1}
Let $\mu$ be the Haar measure of $G$ with $\mu(G) = 1$. Then a sequence $(x_n)$ of elements of $X$ is $\mu$-equidistributed if and only if for any irreducible character $\chi$ of $G$, $\chi \neq 1$, we have
\[
\lim_{n \to \infty} \frac{1}{n} \sum_{i=1}^{n} \chi(x_i) = 0.
\]
\end{corollary}

This follows from Prop.~\ref{prop:I.A.1.2} and the following facts:
$\mu(\chi) = 0$ if $\chi$ is irreducible $\neq 1$, and $\mu(1) = 1$.

\begin{corollary}[2]\label{cor:I.A.1.2}\texorpdfstring{(H. Weyl \cite{weyl-gleichv})}{}
Let $G = \R/\Z$, and let $\mu$ be the normalized Haar measure on $G$. Then $(x_n)$ is $\mu$-equidistributed if and only if for any integer $m \neq 0$ we have
\[
\sum_{n \leq N} e^{2\pi i m x_n} = o(N) \qquad (N \to \infty).
\]
\end{corollary}

For the proof, it suffices to remark that the irreducible characters of $\R/\Z$ are the mappings $x \mapsto e^{2\pi imx}$ ($m \in \Z$).

\subsechead{A.2}{The connection with \textit{L}-functions}

Let $G$ and $X$ be as in Example~2 above: $G$ a compact group and $X$ the space of its conjugacy classes. Let $x_v$, $v \in \Sigma$, be a family of elements of $X$, indexed by a denumerable set $\Sigma$, and let $v \mapsto Nv$ be a function on $\Sigma$ with values in the set of integers $\geq 2$. We make the following hypotheses:

(1) \emph{The infinite product $\prod_{v \in \Sigma} \frac{1}{1 - (Nv)^{-s}}$ converges for every $s \in \C$ with $\mathrm{R}(s) > 1$, and extends to a meromorphic function on $\mathrm{R}(s) \geq 1$ having neither zero nor pole except for a simple pole at $s = 1$.}

(2) \emph{Let $\rho$ be an irreducible representation of $G$, with character $\chi$, and put $L(s, \rho) = \prod_{v \in \Sigma} \frac{1}{\det(1 - \rho(x_v)(Nv)^{-s})}$. Then this product converges for $\mathrm{R}(s) > 1$, and extends to a meromorphic function on $\mathrm{R}(s) \geq 1$ having neither zero nor pole except possibly for $s = 1$.}

The order of $L(s, \rho)$ at $s = 1$ will be denoted by $c_\chi$. Hence, if $L(s, \rho)$ has a pole (resp. a zero) of order $m$ at $s = 1$, one has $c_\chi = m$ (resp. $c_\chi = -m$). Under these assumptions, we have:

\begin{theorem}[1]\label{thm:I.A.2.1}\texorpdfstring{(Wiener--Ikehara)}{}
(a) The number of $v \in \Sigma$ with $Nv \leq n$ is equivalent to $n/\log n$ (as $n \to \infty$).

(b) For any irreducible character $\chi$ of $G$, we have
\[
\sum_{Nv < n} \chi(x_v) = c_\chi \, n/\log n + o(n/\log n) \qquad (n \to \infty).
\]
\end{theorem}

The theorem results, by a standard argument, from the theorem of Wiener--Ikehara, cf.~A.~3 below.

Suppose now that the function $v \mapsto Nv$ has the following property:

(3) There exists a constant $C$ such that, for every $n \in \Z$, the number of $v \in \Sigma$ with $Nv = n$ is $< C$.

One may then arrange the elements of $\Sigma$ as a sequence $(v_i)_{i \geq 1}$ so that $i < j$ implies $Nv_i \leq Nv_j$ (in general, this is possible in many ways). It then makes sense to speak about the equidistribution of the sequence of $x_v$'s; using (3), one shows easily that this does not depend on the chosen ordering of $\Sigma$. Applying Theorem~\ref{thm:I.A.2.1} and Proposition~\ref{prop:I.A.1.2}, we obtain:

\begin{theorem}[2]\label{thm:I.A.2.2}
The elements $x_v$ ($v \in \Sigma$) are equidistributed in $X$ with respect to a measure $\mu$ such that for any irreducible character $\chi$ of $G$ we have $\mu(\chi) = c_\chi$.
\end{theorem}

\begin{corollary}
The elements $x_v$ ($v \in \Sigma$) are equidistributed for the normalized Haar measure of $G$ if and only if $c_\chi = 0$ for every irreducible character $\chi \neq 1$ of $G$, i.e. if and only if the $L$-functions relative to the non-trivial irreducible characters of $G$ are holomorphic and non-zero at $s = 1$.
\end{corollary}

\begin{remarks}
\begin{nlist}
\item Let $G$ be the Galois group of a finite Galois extension $L/K$ of the number field $K$, let $\Sigma$ be the set of unramified places of $K$, let $x_v$ be the Frobenius conjugacy class defined by $v \in \Sigma$, and let $Nv$ be the norm of $v$, cf.~\S\ref{sec:I.2.1}. Properties (1), (2), (3) are satisfied with $c_\chi = 0$ for all irreducible $\chi \neq 1$. This is trivial for (3). For (1), one remarks that $L(s,1)$ is the zeta function of $K$ (up to a finite number of terms), hence has a simple pole at $s = 1$ and is holomorphic on the rest of the line $\mathrm{R}(s) = 1$, cf.~for instance Lang \cite{lang-algnum}, chap.~VII; for a proof of (2) cf.~Artin \cite{artin-collected}, p.~121. Hence Theorem~\ref{thm:I.A.2.2} gives the equidistribution of the Frobenius elements, i.e. the Cebotarev density theorem, cf.~\S\ref{sec:I.2.2}.

\item Let $C$ be the id\`ele class group of a number field $K$, and let $\rho$ be a continuous homomorphism of $C$ into a compact abelian Lie group $G$. An easy argument (cf.~chap.~\ref{chap:III}, \S\ref{sec:III.2.2}) shows that $\rho$ is almost everywhere unramified (i.e. if $U_v$ denotes the group of units at $v$, then $\rho(U_v) = 1$ for almost all $v$). Choose $\pi_v \in K^*$ with $v(\pi_v) = 1$. If $\rho$ is unramified at $v$, then $\rho(\pi_v)$ depends only on $v$, and we set $x_v = \rho(\pi_v)$. We make the following assumption:

($*$) \emph{The homomorphism $\rho$ maps the group $C^0$ of id\`eles of volume $1$ onto $G$.}

(Recall that the \emph{volume} of an id\`ele $a = (a_v)$ is defined as the product of the normalized absolute values of its components $a_v$, cf.~Lang \cite{lang-algnum} or Weil \cite{weil-bnt}.)

Then the elements $x_v$ are uniformly distributed in $G$ with respect to the normalized Haar measure. This follows from Theorem~\ref{thm:I.A.2.2} and the fact that the $L$-functions relative to the irreducible characters $\chi$ of $G$ are Hecke $L$-functions with Gr\"ossencharacters; these $L$-functions are holomorphic and non-zero for $\mathrm{R}(s) \geq 1$ if $\chi \neq 1$, see \cite{lang-algnum}, chap.~VII.
\end{nlist}
\end{remarks}

\begin{remarks}
This example (essentially due to Hecke) is given in Lang (loc.~cit., chap.~VIII, \S 5) except that Lang has replaced the condition ($*$') by the condition ``$\rho$ is surjective'', which is insufficient. This led him to affirm that, for example, the sequence $(\log p)$ (and also the sequence $(\log n)$) is uniformly distributed modulo $1$; however, one knows that this sequence is not uniformly distributed for any measure on $\R/\Z$ (cf.~P\'olya--Szeg\"o \cite{polya-szego}, p.~179--180).
\end{remarks}

\begin{remarks}
(Conjectural example.) Let $E$ be an elliptic curve defined over a number field $K$ and let $\Sigma$ be the set of finite places $v$ of $K$ such that $E$ has good reduction at $v$, cf.~\S\ref{sec:I.1.2} and chap.~IV. Let $v \in \Sigma$, let $\ell \neq p_v$ and let $F_v$ be the Frobenius conjugacy class of $v$ in $\Aut(T_\ell(E))$. The eigenvalues of $F_v$ are algebraic numbers; when embedded into $\C$ they give conjugate complex numbers $\pi_v$, $\bar{\pi}_v$ with $|\pi_v| = Nv^{1/2}$. We may write then $\pi_v = (Nv)^{1/2} e^{i\phi_v}$, $\bar{\pi}_v = (Nv)^{1/2} e^{-i\phi_v}$ with $0 \leq \phi_v \leq \pi$.

On the other hand, let $G = \SU(2)$ be the Lie group of $2 \times 2$ unitary matrices with determinant $1$. Any element of the space $X$ of conjugacy classes of $G$ contains a unique matrix of the form
\[
\begin{pmatrix}
e^{i\phi} & 0 \\ 0 & e^{-i\phi}
\end{pmatrix}, \qquad 0 \leq \phi \leq \pi.
\]
The image in $X$ of the Haar measure of $G$ is known to be $\frac{2}{\pi} \sin^2 \phi \, d\phi$. The irreducible representations of $G$ are the $m$-th symmetric powers $\rho_m$ of the natural representation $\rho_1$ of degree $2$. Take now for $x_v$ the element of $X$ corresponding to the angle $\phi = \phi_v$ defined above. The corresponding $L$-function, relative to $\rho_m$, is:
\[
L_{\rho_m}(s) = \prod_{v} \prod_{a=0}^{m} \frac{1}{1 - \pi_v^{m-a} \bar{\pi}_v^{\,a} (Nv)^{-s}} = L^{\mathfrak{m}}_m(s - m/2).
\]
The function $L^{\mathfrak{m}}_m$ has been considered by Tate \cite{tate-cycles}. He conjectures that $L^{\mathfrak{m}}_m$, for $m \geq 1$, is holomorphic and non-zero for $\mathrm{R}(s) \geq 1 + m/2$, provided that $E$ has no complex multiplication. Granting this conjecture, the corollary to Theorem~\ref{thm:I.A.2.2} would yield the uniform distribution of the $x_v$'s, or, equivalently, that the angles $\phi_v$ of the Frobenius elements are uniformly distributed in $[0,\pi]$ with respect to the measure $\frac{2}{\pi} \sin^2 \phi \, d\phi$ (``conjecture of Sato--Tate'').

One can expect analogous results to be true for other $\ell$-adic representations.
\end{remarks}

\subsechead{A.3}{Proof of theorem~\ref{thm:I.A.2.1}}

The logarithmic derivative of $L$ is
\[
L'/L = \sum_{v,m \geq 1} \frac{\chi(x_v^{m}) \log(Nv)}{(Nv)^{ms}},
\]
where $x_v^{m}$ is the conjugacy class consisting of the $m$-th powers of elements in the class $x_v$. One sees this by writing $L$ as the product $\prod_{i,v} \frac{1}{1 - \lambda_v^{(i)} (Nv)^{-s}}$, where the $\lambda_v^{(i)}$ are the eigenvalues of $x_v$ in the given representation.

Now the series $\sum_{v,\,m \geq 2} \frac{\log(Nv)}{|(Nv)^{ms}|}$ converges for $\mathrm{R}(s) > 1/2$. Indeed, it suffices to show that $\sum_v \frac{\log(Nv)}{(Nv)^\sigma} < \infty$ if $\sigma > 1$; but this series is majorized by $(\mathrm{Constant}) \times \sum_v \frac{1}{(Nv)^{\sigma + \varepsilon}}$ ($\varepsilon > 0$). On the other hand, the convergence for $\sigma > 1$ of the product $\prod_v \frac{1}{1 - (Nv)^{-\sigma}}$ shows that $\sum_v \frac{1}{(Nv)^\sigma} < \infty$ for $\sigma > 1$; hence our assertion. One can therefore write
\[
L'/L = -\sum_v \frac{\chi(x_v) \log(Nv)}{(Nv)^s} + \phi(s),
\]
where $\phi(s)$ is holomorphic for $\mathrm{R}(s) > \frac{1}{2}$. Moreover, by hypothesis, $L'/L$ can be extended to a meromorphic function on $\mathrm{R}(s) \geq 1$ which is holomorphic except possibly for a simple pole at $s = 1$ with residue $-c_\chi$. One may then apply the Wiener--Ikehara theorem (cf.~\cite{lang-algnum}, p.~123):

\begin{theorem}\label{thm:wiener-ikehara}
Let $\mathbb{F}(s) = \sum a_n / n^s$ be a Dirichlet series with complex coefficients. Suppose there exists a Dirichlet series $\mathbb{F}^+(s) = \sum a_n^+ / n^s$ with positive real coefficients such that:
\begin{nlist}
\item $|a_n| \leq a_n^+$ for all $n$;
\item The series $\mathbb{F}^+$ converges for $\mathrm{R}(s) > 1$;
\item The function $\mathbb{F}^+$ (resp. $\mathbb{F}$) can be extended to a meromorphic function on $\mathrm{R}(s) \geq 1$ having no poles except (resp. except possibly) for a simple pole at $s = 1$ with residue $c_+ > 0$ (resp. $c$).
\end{nlist}
Then $\sum_{m<n} a_n = cn + o(n)$ ($n \to \infty$), (where $c = 0$ if $\mathbb{F}$ is holomorphic at $s = 1$).
\end{theorem}

One applies this theorem to
\[
\mathbb{F}(s) = -\sum_v \frac{\chi(x_v) \log(Nv)}{(Nv)^s},
\]
and we take for $\mathbb{F}^+$ the series
\[
d \sum_v \frac{\log(Nv)}{(Nv)^s},
\]
where $d$ is the degree of the given representation $\rho$; this is possible since $\chi(x_v)$ is a sum of $d$ complex numbers of absolute value $1$, hence $|\chi(x_v)| \leq d$; moreover, the series $\sum_v \frac{\log(Nv)}{(Nv)^s}$ differs from the logarithmic derivative of $\prod_v \frac{1}{1 - (Nv)^{-s}}$ by a function which is holomorphic for $\mathrm{R}(s) > 1/2$ as we saw above. Hence by the Wiener--Ikehara theorem we have
\[
\sum_{Nv < n} \chi(x_v) \log(Nv) = c_x n + o(n) \qquad (n \to \infty).
\]
Consequently, by the Abel summation trick (cf.~\cite{lang-algnum}, p.~124, prop.~1),
\[
\sum_{Nv < n} \chi(x_v) = c_\chi \, n/\log n + o(n/\log n) \qquad (n \to \infty),
\]
and in particular,
\[
\sum_{Nv < n} 1 = n/\log n + o(n/\log n) \qquad (n \to \infty).
\]
Hence
\[
\bigl(\sum_{Nv < n} \chi(x_v)\bigr) / \bigl(\sum_{Nv < n} 1\bigr) \longrightarrow c_\chi \qquad \text{as } n \to \infty,
\]
and we may apply Proposition~\ref{prop:I.A.1.2} to conclude the proof. $\square$
